\documentclass[12pt, a4paper]{article}

\usepackage[a4paper, margin=1in]{geometry}

\usepackage[utf8]{inputenc}
\usepackage[mathscr]{euscript}
\usepackage{amssymb,amsmath,amsthm,amsfonts}
\usepackage[shortlabels]{enumitem}
\usepackage{multicol,multirow}
\usepackage{lipsum}
\usepackage{balance}
\usepackage{calc}
\usepackage[colorlinks=true,citecolor=blue,linkcolor=blue]{hyperref}
\usepackage{import}
\usepackage{xifthen}
\usepackage{pdfpages}
\usepackage{transparent}
\usepackage{listings}

\newcommand{\incfig}[2][1.0]{
    \def\svgwidth{#1\columnwidth}
    \import{./figures/}{#2.pdf_tex}
}

\newcommand{\incimg}[2][1.0]{
  \includegraphics[width=#1\columnwidth]{./figures/#2}
}

\newlist{enumproof}{enumerate}{4}
\setlist[enumproof,1]{label=\arabic*., parsep=1em}
\setlist[enumproof,2]{label=\arabic{enumproofi}.\arabic*., parsep=1em}
\setlist[enumproof,3]{label=\arabic{enumproofi}.\arabic{enumproofii}.\arabic*., parsep=1em}
\setlist[enumproof,4]{label=\arabic{enumproofi}.\arabic{enumproofii}.\arabic{enumproofiii}.\arabic*., parsep=1em}

\renewcommand{\qedsymbol}{\ensuremath{\blacksquare}}

\lstdefinestyle{mystyle}{
  language=Python, % Set the language to Python
  commentstyle=\color{codegray}\itshape, % Python comments (#)
  keywordstyle=\color{orange}\bfseries,  % Python reserved keywords
  stringstyle=\color{mauve},             % Strings
  basicstyle={\ttfamily\footnotesize},   % Monospace font
  breakatwhitespace=true,         
  breaklines=true,                 
  captionpos=b,                    
  keepspaces=true,                 
  numbers=none,                    
  tabsize=2,
  showstringspaces=false,
  % Extra keyword groups for Python
  morekeywords=[2]{self, cls, __init__, __str__, __repr__}, 
  keywordstyle=[2]\color{identblue},     % Special identifiers / methods
  morekeywords=[3]{int, float, bool, str, list, dict, set, tuple, object, None, True, False},
  keywordstyle=[3]\color{codegreen},     % Built-in types
  morekeywords=[4]{print, input, open, range, len, type, isinstance, super, enumerate, zip, map, filter, sum, min, max, sorted},
  keywordstyle=[4]\color{cyan},          % Built-in functions
}

% Things Lie
\newcommand{\kb}{\mathfrak b}
\newcommand{\kg}{\mathfrak g}
\newcommand{\kh}{\mathfrak h}
\newcommand{\kn}{\mathfrak n}
\newcommand{\ku}{\mathfrak u}
\newcommand{\kz}{\mathfrak z}
\DeclareMathOperator{\Ext}{Ext} % Ext functor
\DeclareMathOperator{\Tor}{Tor} % Tor functor
\newcommand{\gl}{\opname{\mathfrak{gl}}} % frak gl group
\renewcommand{\sl}{\opname{\mathfrak{sl}}} % frak sl group chktex 6

% More script letters etc.
\newcommand{\SA}{\mathcal A}
\newcommand{\SB}{\mathcal B}
\newcommand{\SC}{\mathcal C}
\newcommand{\SF}{\mathcal F}
\newcommand{\SG}{\mathcal G}
\newcommand{\SH}{\mathcal H}
\newcommand{\OO}{\mathcal O}

\newcommand{\SCA}{\mathscr A}
\newcommand{\SCB}{\mathscr B}
\newcommand{\SCC}{\mathscr C}
\newcommand{\SCD}{\mathscr D}
\newcommand{\SCE}{\mathscr E}
\newcommand{\SCF}{\mathscr F}
\newcommand{\SCG}{\mathscr G}
\newcommand{\SCH}{\mathscr H}

% Mathfrak primes
\newcommand{\km}{\mathfrak m}
\newcommand{\kp}{\mathfrak p}
\newcommand{\kq}{\mathfrak q}

% number sets
\newcommand{\RR}[1][]{\ensuremath{\ifstrempty{#1}{\mathbb{R}}{\mathbb{R}^{#1}}}}
\newcommand{\NN}[1][]{\ensuremath{\ifstrempty{#1}{\mathbb{N}}{\mathbb{N}^{#1}}}}
\newcommand{\ZZ}[1][]{\ensuremath{\ifstrempty{#1}{\mathbb{Z}}{\mathbb{Z}^{#1}}}}
\newcommand{\QQ}[1][]{\ensuremath{\ifstrempty{#1}{\mathbb{Q}}{\mathbb{Q}^{#1}}}}
\newcommand{\CC}[1][]{\ensuremath{\ifstrempty{#1}{\mathbb{C}}{\mathbb{C}^{#1}}}}
\newcommand{\PP}[1][]{\ensuremath{\ifstrempty{#1}{\mathbb{P}}{\mathbb{P}^{#1}}}}
\newcommand{\HH}[1][]{\ensuremath{\ifstrempty{#1}{\mathbb{H}}{\mathbb{H}^{#1}}}}
\newcommand{\FF}[1][]{\ensuremath{\ifstrempty{#1}{\mathbb{F}}{\mathbb{F}^{#1}}}}
% expected value
\newcommand{\EE}{\ensuremath{\mathbb{E}}}
\newcommand{\charin}{\text{ char }}
\DeclareMathOperator{\sign}{sign}
\DeclareMathOperator{\Aut}{Aut}
\DeclareMathOperator{\Inn}{Inn}
\DeclareMathOperator{\Syl}{Syl}
\DeclareMathOperator{\Gal}{Gal}
\DeclareMathOperator{\GL}{GL} % General linear group
\DeclareMathOperator{\SL}{SL} % Special linear group

%---------------------------------------
% BlackBoard Math Fonts :-
%---------------------------------------

%Captital Letters
\newcommand{\bbA}{\mathbb{A}}	\newcommand{\bbB}{\mathbb{B}}
\newcommand{\bbC}{\mathbb{C}}	\newcommand{\bbD}{\mathbb{D}}
\newcommand{\bbE}{\mathbb{E}}	\newcommand{\bbF}{\mathbb{F}}
\newcommand{\bbG}{\mathbb{G}}	\newcommand{\bbH}{\mathbb{H}}
\newcommand{\bbI}{\mathbb{I}}	\newcommand{\bbJ}{\mathbb{J}}
\newcommand{\bbK}{\mathbb{K}}	\newcommand{\bbL}{\mathbb{L}}
\newcommand{\bbM}{\mathbb{M}}	\newcommand{\bbN}{\mathbb{N}}
\newcommand{\bbO}{\mathbb{O}}	\newcommand{\bbP}{\mathbb{P}}
\newcommand{\bbQ}{\mathbb{Q}}	\newcommand{\bbR}{\mathbb{R}}
\newcommand{\bbS}{\mathbb{S}}	\newcommand{\bbT}{\mathbb{T}}
\newcommand{\bbU}{\mathbb{U}}	\newcommand{\bbV}{\mathbb{V}}
\newcommand{\bbW}{\mathbb{W}}	\newcommand{\bbX}{\mathbb{X}}
\newcommand{\bbY}{\mathbb{Y}}	\newcommand{\bbZ}{\mathbb{Z}}

%---------------------------------------
% MathCal Fonts :-
%---------------------------------------

%Captital Letters
\newcommand{\mcA}{\mathcal{A}}	\newcommand{\mcB}{\mathcal{B}}
\newcommand{\mcC}{\mathcal{C}}	\newcommand{\mcD}{\mathcal{D}}
\newcommand{\mcE}{\mathcal{E}}	\newcommand{\mcF}{\mathcal{F}}
\newcommand{\mcG}{\mathcal{G}}	\newcommand{\mcH}{\mathcal{H}}
\newcommand{\mcI}{\mathcal{I}}	\newcommand{\mcJ}{\mathcal{J}}
\newcommand{\mcK}{\mathcal{K}}	\newcommand{\mcL}{\mathcal{L}}
\newcommand{\mcM}{\mathcal{M}}	\newcommand{\mcN}{\mathcal{N}}
\newcommand{\mcO}{\mathcal{O}}	\newcommand{\mcP}{\mathcal{P}}
\newcommand{\mcQ}{\mathcal{Q}}	\newcommand{\mcR}{\mathcal{R}}
\newcommand{\mcS}{\mathcal{S}}	\newcommand{\mcT}{\mathcal{T}}
\newcommand{\mcU}{\mathcal{U}}	\newcommand{\mcV}{\mathcal{V}}
\newcommand{\mcW}{\mathcal{W}}	\newcommand{\mcX}{\mathcal{X}}
\newcommand{\mcY}{\mathcal{Y}}	\newcommand{\mcZ}{\mathcal{Z}}


%---------------------------------------
% Bold Math Fonts :-
%---------------------------------------

%Captital Letters
\newcommand{\bmA}{\boldsymbol{A}}	\newcommand{\bmB}{\boldsymbol{B}}
\newcommand{\bmC}{\boldsymbol{C}}	\newcommand{\bmD}{\boldsymbol{D}}
\newcommand{\bmE}{\boldsymbol{E}}	\newcommand{\bmF}{\boldsymbol{F}}
\newcommand{\bmG}{\boldsymbol{G}}	\newcommand{\bmH}{\boldsymbol{H}}
\newcommand{\bmI}{\boldsymbol{I}}	\newcommand{\bmJ}{\boldsymbol{J}}
\newcommand{\bmK}{\boldsymbol{K}}	\newcommand{\bmL}{\boldsymbol{L}}
\newcommand{\bmM}{\boldsymbol{M}}	\newcommand{\bmN}{\boldsymbol{N}}
\newcommand{\bmO}{\boldsymbol{O}}	\newcommand{\bmP}{\boldsymbol{P}}
\newcommand{\bmQ}{\boldsymbol{Q}}	\newcommand{\bmR}{\boldsymbol{R}}
\newcommand{\bmS}{\boldsymbol{S}}	\newcommand{\bmT}{\boldsymbol{T}}
\newcommand{\bmU}{\boldsymbol{U}}	\newcommand{\bmV}{\boldsymbol{V}}
\newcommand{\bmW}{\boldsymbol{W}}	\newcommand{\bmX}{\boldsymbol{X}}
\newcommand{\bmY}{\boldsymbol{Y}}	\newcommand{\bmZ}{\boldsymbol{Z}}
%Small Letters
\newcommand{\bma}{\boldsymbol{a}}	\newcommand{\bmb}{\boldsymbol{b}}
\newcommand{\bmc}{\boldsymbol{c}}	\newcommand{\bmd}{\boldsymbol{d}}
\newcommand{\bme}{\boldsymbol{e}}	\newcommand{\bmf}{\boldsymbol{f}}
\newcommand{\bmg}{\boldsymbol{g}}	\newcommand{\bmh}{\boldsymbol{h}}
\newcommand{\bmi}{\boldsymbol{i}}	\newcommand{\bmj}{\boldsymbol{j}}
\newcommand{\bmk}{\boldsymbol{k}}	\newcommand{\bml}{\boldsymbol{l}}
\newcommand{\bmm}{\boldsymbol{m}}	\newcommand{\bmn}{\boldsymbol{n}}
\newcommand{\bmo}{\boldsymbol{o}}	\newcommand{\bmp}{\boldsymbol{p}}
\newcommand{\bmq}{\boldsymbol{q}}	\newcommand{\bmr}{\boldsymbol{r}}
\newcommand{\bms}{\boldsymbol{s}}	\newcommand{\bmt}{\boldsymbol{t}}
\newcommand{\bmu}{\boldsymbol{u}}	\newcommand{\bmv}{\boldsymbol{v}}
\newcommand{\bmw}{\boldsymbol{w}}	\newcommand{\bmx}{\boldsymbol{x}}
\newcommand{\bmy}{\boldsymbol{y}}	\newcommand{\bmz}{\boldsymbol{z}}

%---------------------------------------
% Scr Math Fonts :-
%---------------------------------------

\newcommand{\sA}{{\mathscr{A}}}   \newcommand{\sB}{{\mathscr{B}}}
\newcommand{\sC}{{\mathscr{C}}}   \newcommand{\sD}{{\mathscr{D}}}
\newcommand{\sE}{{\mathscr{E}}}   \newcommand{\sF}{{\mathscr{F}}}
\newcommand{\sG}{{\mathscr{G}}}   \newcommand{\sH}{{\mathscr{H}}}
\newcommand{\sI}{{\mathscr{I}}}   \newcommand{\sJ}{{\mathscr{J}}}
\newcommand{\sK}{{\mathscr{K}}}   \newcommand{\sL}{{\mathscr{L}}}
\newcommand{\sM}{{\mathscr{M}}}   \newcommand{\sN}{{\mathscr{N}}}
\newcommand{\sO}{{\mathscr{O}}}   \newcommand{\sP}{{\mathscr{P}}}
\newcommand{\sQ}{{\mathscr{Q}}}   \newcommand{\sR}{{\mathscr{R}}}
\newcommand{\sS}{{\mathscr{S}}}   \newcommand{\sT}{{\mathscr{T}}}
\newcommand{\sU}{{\mathscr{U}}}   \newcommand{\sV}{{\mathscr{V}}}
\newcommand{\sW}{{\mathscr{W}}}   \newcommand{\sX}{{\mathscr{X}}}
\newcommand{\sY}{{\mathscr{Y}}}   \newcommand{\sZ}{{\mathscr{Z}}}


%---------------------------------------
% Math Fraktur Font
%---------------------------------------

%Captital Letters
\newcommand{\mfA}{\mathfrak{A}}	\newcommand{\mfB}{\mathfrak{B}}
\newcommand{\mfC}{\mathfrak{C}}	\newcommand{\mfD}{\mathfrak{D}}
\newcommand{\mfE}{\mathfrak{E}}	\newcommand{\mfF}{\mathfrak{F}}
\newcommand{\mfG}{\mathfrak{G}}	\newcommand{\mfH}{\mathfrak{H}}
\newcommand{\mfI}{\mathfrak{I}}	\newcommand{\mfJ}{\mathfrak{J}}
\newcommand{\mfK}{\mathfrak{K}}	\newcommand{\mfL}{\mathfrak{L}}
\newcommand{\mfM}{\mathfrak{M}}	\newcommand{\mfN}{\mathfrak{N}}
\newcommand{\mfO}{\mathfrak{O}}	\newcommand{\mfP}{\mathfrak{P}}
\newcommand{\mfQ}{\mathfrak{Q}}	\newcommand{\mfR}{\mathfrak{R}}
\newcommand{\mfS}{\mathfrak{S}}	\newcommand{\mfT}{\mathfrak{T}}
\newcommand{\mfU}{\mathfrak{U}}	\newcommand{\mfV}{\mathfrak{V}}
\newcommand{\mfW}{\mathfrak{W}}	\newcommand{\mfX}{\mathfrak{X}}
\newcommand{\mfY}{\mathfrak{Y}}	\newcommand{\mfZ}{\mathfrak{Z}}
%Small Letters
\newcommand{\mfa}{\mathfrak{a}}	\newcommand{\mfb}{\mathfrak{b}}
\newcommand{\mfc}{\mathfrak{c}}	\newcommand{\mfd}{\mathfrak{d}}
\newcommand{\mfe}{\mathfrak{e}}	\newcommand{\mff}{\mathfrak{f}}
\newcommand{\mfg}{\mathfrak{g}}	\newcommand{\mfh}{\mathfrak{h}}
\newcommand{\mfi}{\mathfrak{i}}	\newcommand{\mfj}{\mathfrak{j}}
\newcommand{\mfk}{\mathfrak{k}}	\newcommand{\mfl}{\mathfrak{l}}
\newcommand{\mfm}{\mathfrak{m}}	\newcommand{\mfn}{\mathfrak{n}}
\newcommand{\mfo}{\mathfrak{o}}	\newcommand{\mfp}{\mathfrak{p}}
\newcommand{\mfq}{\mathfrak{q}}	\newcommand{\mfr}{\mathfrak{r}}
\newcommand{\mfs}{\mathfrak{s}}	\newcommand{\mft}{\mathfrak{t}}
\newcommand{\mfu}{\mathfrak{u}}	\newcommand{\mfv}{\mathfrak{v}}
\newcommand{\mfw}{\mathfrak{w}}	\newcommand{\mfx}{\mathfrak{x}}
\newcommand{\mfy}{\mathfrak{y}}	\newcommand{\mfz}{\mathfrak{z}}


\newcommand{\mytitle}{FMA\_3TC13\_TP Colle 1}
\newcommand{\myauthor}{github/omgeta}
\newcommand{\mydate}{AU 25-26 S1}

\begin{document}
\raggedright
\footnotesize
\begin{center}
{\normalsize{\textbf{\mytitle}}} \\
{\footnotesize{\mydate\hspace{2pt}\textemdash\hspace{2pt}\myauthor}}
\end{center}
\setlist{topsep=-1em, itemsep=-1em, parsep=2em}
%%%%%%%%%%%%%%%%%%%%%%%%%%%%%%%%%%%%%%%%%%%%%%%%%%%%%%
%                      Begin                         %
%%%%%%%%%%%%%%%%%%%%%%%%%%%%%%%%%%%%%%%%%%%%%%%%%%%%%%
\begin{enumerate}
  \item $A = \left(\begin{array}{ccccc} 1 & 1 & 2 & 1 & 1\\ 2 & 1 & 1 & 1 & 1\\ 1 & 1 & 1 & 2 & 1\\ 2 & 1 & 1 & 1 & 1\\ 1 & 1 & 1 & 1 & 2 \end{array}\right) \xrightarrow{RREF} \left(\begin{array}{ccccc} 1 & 0 & 0 & 0 & -1\\ 0 & 1 & 0 & 0 & 5\\ 0 & 0 & 1 & 0 & -1\\ 0 & 0 & 0 & 1 & -1\\ 0 & 0 & 0 & 0 & 0 \end{array}\right)$

    Alors avec 4 pivots, on a $\boxed{rg(A) = 4}$

    Par la réduction, on a la base échelonnée reduite de l'espace des lignes:
    \[
    \boxed{
    \left\{
    (1,0,0,0,-1),\,
    (0,1,0,0,5),\,
    (0,0,1,0,-1),\,
    (0,0,0,1,-1)
    \right\}.
    }
    \]

    Puis, par l'observation, on a la base échelonnée réduite de l'espace des colonnes:
    \[
    \boxed{
    \left\{
    \left(\begin{array}{c}1\\0\\0\\0\\0\end{array}\right),
    \left(\begin{array}{c}0\\1\\0\\1\\0\end{array}\right),
    \left(\begin{array}{c}0\\0\\1\\0\\0\end{array}\right),
    \left(\begin{array}{c}0\\0\\0\\0\\1\end{array}\right)
    \right\}.
    }
    \]
    \pagebreak
  \item 
    On peut écrire
    \[
    B=AC
    \]
    avec
    \[
    C=
    \left(\begin{array}{ccccc}
    0&1&1&\cdots&1\\
    1&0&1&\cdots&1\\
    1&1&0&\cdots&1\\
    \vdots&\vdots&\vdots&\ddots&\vdots\\
    1&1&1&\cdots&0
    \end{array}\right).
    \]

    Donc, par multiplicativité du déterminant,
    \[
    \det(B)=\det(A)\det(C).
    \]

    Calculons $\det(C)$. On effectue
    \[
    C_1\leftarrow C_1+C_2+\cdots+C_n.
    \]
    La première colonne devient
    \[
    (n-1)
    \left(\begin{array}{c}
    1\\
    1\\
    \vdots\\
    1
    \end{array}\right).
    \]

    On factorise $(n-1)$ puis, pour $j=2,\ldots,n$, on effectue
    \[
    C_j\leftarrow C_j-C_1.
    \]

    On obtient alors
    \[
    \det(C)
    =
    (n-1)
    \left|
    \begin{array}{ccccc}
    1&0&0&\cdots&0\\
    1&-1&0&\cdots&0\\
    1&0&-1&\cdots&0\\
    \vdots&\vdots&\vdots&\ddots&\vdots\\
    1&0&0&\cdots&-1
    \end{array}
    \right|.
    \]

    La matrice est triangulaire, donc
    \[
    \det(C)=(n-1)(-1)^{n-1}.
    \]

    Ainsi
    \[
    \boxed{\det(B)=(-1)^{n-1}(n-1)\det(A).}
    \]

    \pagebreak 

  \item 
    En développant par rapport à la première ligne, puis à la première colonne du
mineur obtenu, on trouve la relation de récurrence
    \[
    \boxed{D_n(x)=2\cos x\,D_{n-1}(x)-D_{n-2}(x)}.
    \]

    De plus,
    \[
    D_0(x)=1,
    \qquad
    D_1(x)=2\cos x.
    \]

    Pour $\sin x\neq 0$, on utilise l'identité
    \[
    \sin((n+1)x)
    =
    2\cos x\,\sin(nx)-\sin((n-1)x).
    \]

    La suite
    \[
    u_n=\frac{\sin((n+1)x)}{\sin x}
    \]
    vérifie donc la même relation de récurrence :
    \[
    u_n=2\cos x\,u_{n-1}-u_{n-2}.
    \]

    De plus,
    \[
    u_0=1=D_0,
    \qquad
    u_1=\frac{\sin(2x)}{\sin x}=2\cos x=D_1.
    \]

    Par récurrence,
    \[
    \boxed{
    D_n(x)=\frac{\sin((n+1)x)}{\sin x}
    }
    \qquad \text{si } \sin x\neq 0.
    \]

    Si $x=k\pi$, alors $\cos x=(-1)^k$ et la relation de récurrence donne
    \[
    \boxed{D_n(k\pi)=(-1)^{kn}(n+1)}.
    \]

    Ainsi,
    \[
    \boxed{
    D_n(x)=
    \begin{cases}
    \dfrac{\sin((n+1)x)}{\sin x},
    & x\notin \pi\mathbb Z,\\[6pt]
    (-1)^{kn}(n+1),
    & x=k\pi.
    \end{cases}
    }
    \]
  \pagebreak 

\item 
  \begin{enumerate}[\arabic*.]
    \item 
      \begin{enumproof}
      \item ($\Rightarrow$) 
        \begin{enumproof}
        \item Si $F \subseteq G$, $F \cup G = G$ qui est déjà sous-espace de $E$
        \item Si $G \subseteq F$, $F \cup G = F$ qui est déjà sous-espace de $E$
        \item $\therefore F \cup G$ est un sous-espace vectoriel de $E$
        \end{enumproof}
      \item ($\Leftarrow$) 
        \begin{enumproof}
        \item Supposons la converse, $F\cup G$ est un sous-espace mais $F\not\subseteq G$ et $G\not\subseteq F$
        \item $\exists x \in F \setminus G,\; y \in G\setminus F$
        \item $(x+y)\in F\cup G$\hfill(Definition de sous-espace)
        \item Si $(x+y)\in F$, $(x+y)-x= y \in F$
        \item Si $(x+y)\in G$, $(x+y)-y= x \in G$
        \item $\therefore$ par contradiction, $F\subseteq G \lor G \subseteq F$
        \end{enumproof}
      \end{enumproof}

    \item \begin{enumproof}
      \item ($\subseteq$)
        \begin{enumproof}
          \item Soit $x\in F\cap(G+H)$
          \item Alors $x\in F$ et $\exists g\in G,\; h\in H$ tels que $x=g+h$
          \item Comme $G\subseteq F$, $g\in F$
          \item Donc $h=x-g\in F$, ainsi $h\in F\cap H$
          \item $\therefore x\in G+(F\cap H)$
        \end{enumproof}

      \item ($\supseteq$)
        \begin{enumproof}
          \item Soit $x\in G+(F\cap H)$
          \item Alors $\exists g\in G,\; h\in F\cap H$ tels que $x=g+h$
          \item Comme $G\subseteq F$, $g\in F$, et $h\in F$
          \item Donc $x\in F$
          \item De plus, $g\in G$ et $h\in H$, donc $x\in G+H$
          \item $\therefore x\in F\cap(G+H)$
        \end{enumproof}

      \item $\therefore F\cap(G+H)=G+(F\cap H)$
    \end{enumproof}
  \end{enumerate}

  \item Posons $\left(\begin{array}{ccccc} 1 & 1 & 3 & 10 & 1\\ 1 & 2 & 1 & 4 & 7\\ 1 & 3 & 4 & 13 & 8\\ 1 & 4 & 2 & 7 & 14 \end{array}\right) \xrightarrow{RREF} \left(\begin{array}{ccccc} 1 & 0 & 0 & 1 & 0\\ 0 & 1 & 0 & 0 & 4\\ 0 & 0 & 1 & 3 & -1\\ 0 & 0 & 0 & 0 & 0 \end{array}\right)$ 

    Donc, $(v_1, v_2, v_3)$ est une base de $F$ avec $dim(F) = 3$

    Puis, réduire le système homogène $\left(\begin{array}{ccccc} 1 & 1 & 1 & 1 & 0\\ 1 & 2 & 3 & 4 & 0\\ 3 & 1 & 4 & 2 & 0 \end{array}\right) \xrightarrow{RREF} \left(\begin{array}{ccccc} 1 & 0 & 0 & -1 & 0\\ 0 & 1 & 0 & 1 & 0\\ 0 & 0 & 1 & 1 & 0 \end{array}\right)$

    Ce qui donne la solution $(a, b, c, d) = t(1, -1, -1, 1)$ et si on prend $t=1$, on obtient $\boxed{x_1 - x_2 - x_3 + x_4 = 0}$

    \pagebreak
  \item On développe :
    \[
    P_0=1-\frac32X+\frac12X^2,\qquad
    P_1=-\frac12X+\frac12X^2,
    \]
    \[
    P_2=-4X+2X^2,\qquad
    P_3=1-\frac43X+\frac13X^2.
    \]
    \begin{enumerate}[\arabic*.]
      \item $\left(\begin{array}{ccc} 1 & 0 & 0\\ -\frac{3}{2} & -\frac{1}{2} & -4\\ \frac{1}{2} & \frac{1}{2} & 2 \end{array}\right) \xrightarrow{RREF} I_3$ alors chaque colonne est indép. donc $\mcB$ est une base de $E$

      \item $P_{\mcC, \mcB} = \left(\begin{array}{ccc} 1 & 0 & 0\\ -\frac{3}{2} & -\frac{1}{2} & -4\\ \frac{1}{2} & \frac{1}{2} & 2 \end{array}\right)$ alors $P_{\mcB, \mcC} = P_{\mcC, \mcB}^{-1} = \left(\begin{array}{ccc} 1 & 0 & 0\\ 1 & 2 & 4\\ -\frac{1}{2} & -\frac{1}{2} & -\frac{1}{2} \end{array}\right)$

      \item $P_0, P_1$ sont indép, donc $dim F = 2$\\ 
        $P_2, P_3$ ne sont pas proportional donc $dim G = 2$\\
        $(P_0, P_1, P_2)$ est une base de $E$ avec $dim E = 3$ alors $dim F+G = 3$\\
        Notons $P_3 = P_0 - \frac{1}{3}P_1$ pour montrer que $P_3 \in F$ et $P_3 \in G$ alors $dim F \cap G = 1$

      \item Verifions avec les valeurs trouvé dans la partie 3
    \end{enumerate}
\end{enumerate}
%%%%%%%%%%%%%%%%%%%%%%%%%%%%%%%%%%%%%%%%%%%%%%%%%%%%%%
%                       End                          %
%%%%%%%%%%%%%%%%%%%%%%%%%%%%%%%%%%%%%%%%%%%%%%%%%%%%%%

\end{document}
